\section{Trabajos relacionados}
\label{sec:trabajos-relacionados}

\subsection{Bases de datos distribuidas y consistencia}

CockroachDB es una base de datos SQL distribuida diseñada para tolerar la pérdida de nodos y
regiones enteras sin intervención manual, replicando cada rango de datos mediante el protocolo de
consenso Raft~\cite{ongaro2014raft} y garantizando aislamiento \texttt{SERIALIZABLE} de extremo a
extremo mediante un protocolo de \emph{timestamps} y reintentos transaccionales
transparentes~\cite{taft2020cockroachdb}. Esta decisión de diseño —consistencia fuerte por
defecto, en vez de consistencia eventual configurable— la sitúa, en el espacio de \emph{trade-offs}
descrito por el teorema CAP~\cite{brewer2012cap} y su refinamiento
PACELC~\cite{abadi2012consistency}, del lado de la consistencia incluso a costa de latencia
adicional en el caso de conflicto, algo directamente observable en las pruebas de integración de
la Sección~\ref{sec:cluster-verificacion}: dos transacciones concurrentes sobre la misma fila no
producen un resultado silenciosamente inconsistente, sino un error de conflicto explícito
(\texttt{SQLSTATE 40001}) que la aplicación debe reintentar. Esta elección es consistente con el
dominio del problema: un sistema de tickets de soporte no puede permitirse que dos actualizaciones
concurrentes sobre el mismo ticket (por ejemplo, un técnico cerrándolo mientras un administrador lo
reasigna) se resuelvan de forma silenciosamente incorrecta.

La fragmentación horizontal aplicada en este trabajo se apoya en el marco teórico clásico de
Özsu y Valduriez~\cite{ozsu2011principles}, que exige que toda fragmentación cumpla completitud,
reconstrucción y disyunción (verificado formalmente en la Sección~\ref{sec:diseno-esquema}). A
diferencia de sistemas NoSQL de consistencia eventual descritos en el mismo texto, CockroachDB
mantiene semántica SQL completa y transacciones ACID distribuidas, lo que permitió mantener sin
cambios la capa de persistencia JPA de \texttt{ticket-service} —solo cambió la definición de la
clave primaria y la estrategia de acceso, no el modelo transaccional de la aplicación.

\subsection{Procesamiento paralelo de datos}

El pipeline analítico de este trabajo se construye sobre el modelo de \emph{Resilient Distributed
Datasets} (RDD) de Apache Spark~\cite{zaharia2012rdd}, que generaliza el modelo MapReduce
permitiendo mantener conjuntos de datos intermedios en memoria entre etapas —esencial para un
pipeline de cinco transformaciones encadenadas como el descrito en la
Sección~\ref{sec:pipeline-paralelo}, donde recomputar cada etapa desde disco habría anulado buena
parte de la ganancia de paralelismo—. Salloum et al.~\cite{salloum2016bigdata} documentan
precisamente esta ventaja de Spark frente a MapReduce clásico para cargas de trabajo iterativas y
de aprendizaje automático, que es el caso de la etapa de \emph{clustering} (T5) de este pipeline.
La guía de referencia de Chambers y Zaharia~\cite{chambers2018spark} se usó como referencia práctica
para la API de \texttt{DataFrame} y \texttt{pyspark.ml} empleada en \texttt{transformations.py}.

Para el protocolo de medición de rendimiento, el antecedente directo es el trabajo de
Cooper et al.~\cite{cooper2010ycsb} sobre \emph{benchmarking} sistemático de sistemas de
almacenamiento en la nube (YCSB), cuyo principio metodológico central —repeticiones controladas,
descarte de efectos de arranque en frío, y comparación contra una línea base conocida— se adapta en
este trabajo al protocolo experimental de la Sección~\ref{sec:protocolo-resultados}, aunque aquí el
objeto de medición es un pipeline analítico completo y no operaciones individuales de
lectura/escritura.

\subsection{Metodología experimental}

El diseño del protocolo experimental (variables independientes y dependientes, repeticiones,
descarte de calentamiento/enfriamiento, prueba de normalidad seguida de prueba paramétrica o no
paramétrica según corresponda) sigue las guías de Wohlin et al.~\cite{wohlin2012experimentation}
y los principios de diseño de experimentos de sistemas de Jain~\cite{jain1991art}. El uso del caso
de estudio del propio sistema como unidad de análisis, y la forma de reportar las amenazas a la
validez interna y externa en la Sección~\ref{sec:protocolo-resultados}, siguen las recomendaciones
de Runeson y Höst~\cite{runeson2009guidelines} para investigación de estudios de caso en ingeniería
de software.
